\section{Discussion}
\label{sec:discussion}
\paragraph{How much does interpretability cost?} Very little on this dataset.
The random forest is ahead by 0.008 of AUC, a gap that is consistent across
folds and on the test set but small next to the 0.36 that separates either
model from the baseline. At a fixed recall of 70\% it means 14\% fewer good
payers refused, which may or may not justify a model whose refusals are harder
to explain. This is the setup described in the course for credit: a flexible
model for accuracy and a logistic regression as the explainable reference
\citep{oku2026slides}. The binning experiment suggests where the gap comes
from. Giving the linear model non-linear curves per feature recovered only
0.003 of AUC, so most of the forest's advantage comes from interactions between
features, which a logistic regression only captures if they are created by
hand.

\paragraph{Why regularization did not help.} With 120{,}000 training clients
and 14 coefficients, each coefficient is estimated from thousands of examples
and its variance is already small. A penalty can only add bias, which is what
Figure~\ref{fig:regularization} shows. Ridge and lasso earn their place when
there are few rows, many columns or near-collinear columns
\citep{oku2026slides,james2021islr}, none of which describes this table. The
same figure shows why model selection used the log-loss: a strong ridge
penalty kept the ranking (AUC 0.847) while producing probabilities worse
than a constant, a failure that the AUC cannot see.

\paragraph{Why the data cleaning mattered.} Treating the 96/98 codes as counts,
or simply imputing the missing incomes, would have hidden two signals: 269
clients with a 55\% default rate, and a missing-income group that defaults
less than average. Encoding both as flags, and keeping every learned step
inside the cross-validation folds, is also why the cross-validated and test
scores agree to within 0.007.

\paragraph{Limitations.} The coefficients describe associations conditional on
the other features, not causes; the three late-payment counts are strongly
correlated, so their individual odds ratios share the weight between them, and
impurity-based importances are unsigned and favour features with many distinct
values. The decision threshold follows an illustrative recall target, not the
real costs of each error, and the calibration of the probabilities was not
examined. The dataset has no dates, so a split by time, the realistic test for
a credit model, was not possible; it also contains 609 duplicated rows that may
place the same client on both sides of a split. Only the leaf size of the
forest was tuned, and gradient-boosted trees \citep{chen2016xgboost}, the
usual choice for tabular credit data, were not tried. Finally, age is one of
the strongest predictors, with younger clients scored as riskier. Whether a
lender may use it is a legal and ethical question that the model cannot
answer, and a reminder that a model can learn and reproduce a prejudice
present in its data \citep{oku2026slides}.
